\documentclass[11pt,border=4pt]{standalone}
% Estilo comum dos esquemáticos do relatório (circuitikz).
% Uso: \documentclass[11pt,border=4pt]{standalone} e, em seguida, % Estilo comum dos esquemáticos do relatório (circuitikz).
% Uso: \documentclass[11pt,border=4pt]{standalone} e, em seguida, % Estilo comum dos esquemáticos do relatório (circuitikz).
% Uso: \documentclass[11pt,border=4pt]{standalone} e, em seguida, \input{estilo_circuitos}.
\usepackage[T1]{fontenc}
\usepackage[utf8]{inputenc}
\usepackage[brazilian]{babel}
\usepackage{lmodern}
\usepackage{amsmath}
\usepackage[american]{circuitikz}
\usetikzlibrary{calc,arrows.meta}

% Cores das figuras originais
\definecolor{anot}{RGB}{0,135,200}      % anotações: valores calculados, setas, resistências vistas
\definecolor{sombra}{RGB}{220,242,253}  % região do transistor nos modelos T e pi
\definecolor{term}{RGB}{62,124,148}     % nomes de terminais destacados (B, C, E)

% Fios com 0,4 pt e corpo dos componentes com 0,8 pt, como nas figuras originais
\ctikzset{bipoles/thickness=2}
.
\usepackage[T1]{fontenc}
\usepackage[utf8]{inputenc}
\usepackage[brazilian]{babel}
\usepackage{lmodern}
\usepackage{amsmath}
\usepackage[american]{circuitikz}
\usetikzlibrary{calc,arrows.meta}

% Cores das figuras originais
\definecolor{anot}{RGB}{0,135,200}      % anotações: valores calculados, setas, resistências vistas
\definecolor{sombra}{RGB}{220,242,253}  % região do transistor nos modelos T e pi
\definecolor{term}{RGB}{62,124,148}     % nomes de terminais destacados (B, C, E)

% Fios com 0,4 pt e corpo dos componentes com 0,8 pt, como nas figuras originais
\ctikzset{bipoles/thickness=2}
.
\usepackage[T1]{fontenc}
\usepackage[utf8]{inputenc}
\usepackage[brazilian]{babel}
\usepackage{lmodern}
\usepackage{amsmath}
\usepackage[american]{circuitikz}
\usetikzlibrary{calc,arrows.meta}

% Cores das figuras originais
\definecolor{anot}{RGB}{0,135,200}      % anotações: valores calculados, setas, resistências vistas
\definecolor{sombra}{RGB}{220,242,253}  % região do transistor nos modelos T e pi
\definecolor{term}{RGB}{62,124,148}     % nomes de terminais destacados (B, C, E)

% Fios com 0,4 pt e corpo dos componentes com 0,8 pt, como nas figuras originais
\ctikzset{bipoles/thickness=2}

% Questão 3: meio-circuito diferencial de M2 (a) e seu modelo pi com r_o (b).
% Compilar com pdflatex; o PDF resultante vai para figuras/fig_q3_dif_meio.pdf.
\begin{document}
\begin{circuitikz}
  % ---------------- (a) Meio-circuito diferencial ----------------
  % M2 com o dreno em (-7.6,0); nó de saída 0,4 cm acima do dreno
  \draw (-7.6,0) node[nmos, arrowmos, anchor=D] (M) {};
  \draw (M.D) -- (-7.6,0.4) to[R] (-7.6,2.0) node[ground, rotate=180]{};
  \draw (-7.6,0.4) -- (-6.4,0.4);
  \draw (M.G) -- ++(-0.7,0) coordinate (Ga);
  \draw (M.S) node[ground] (Ta) {};
  \node[circ] at (-7.6,0.4) {};
  \node[ocirc] at (-6.4,0.4) {};
  \node[ocirc] at (Ga) {};
  % rótulos
  \node[anchor=east, inner sep=0pt] at (-7.97,1.2) {$R_D$};
  \node[anchor=west, inner sep=0pt, text=anot] at (-6.184,0.4) {$v_{o2}$};
  \node[anchor=east, inner sep=0pt, text=anot] at ($(Ga)+(-0.216,0)$) {$-v_{id}/2$};
  \node[anchor=west, inner sep=0pt] at ($(M)+(0.25,0)$) {$M_2$};
  \node[anchor=base west, inner sep=0pt, text=anot] at ($(Ta)+(0.37,-0.42)$) {terra virtual};

  % ---------------- (b) Modelo pi ----------------
  % nó D em (0,0); fio de baixo (fonte S) em y = -3,5; porta G na altura do centro da fonte controlada
  \draw (0,0) to[cI] (0,-3.5);
  \draw (0,0) -- (1.4,0) -- (2.7,0) -- (3.6,0);
  \draw (1.4,0) to[R] (1.4,-3.5) -- (0,-3.5);
  \draw (2.7,0) to[R] (2.7,1.7) node[ground, rotate=180]{};
  \draw (-2.2,-3.5) -- (0,-3.5);
  \draw (0.7,-3.5) node[ground]{};
  \foreach \p in {(0,0),(1.4,0),(2.7,0),(0,-3.5)} \node[circ] at \p {};
  \foreach \p in {(3.6,0),(-2.2,-1.75),(-2.2,-3.5)} \node[ocirc] at \p {};
  % rótulos
  \node[anchor=base, inner sep=0pt] at (0,0.27) {$D$};
  \node[anchor=east, inner sep=0pt] at (-0.65,-1.75) {$g_m v_{gs}$};
  \node[anchor=west, inner sep=0pt] at (1.77,-1.75) {$r_o$};
  \node[anchor=west, inner sep=0pt] at (3.07,0.85) {$R_D$};
  \node[anchor=west, inner sep=0pt, text=anot] at (3.816,0) {$v_{o2}$};
  \node[anchor=base, inner sep=0pt] at (-2.2,-1.5) {$G$};
  \node[anchor=base, inner sep=0pt] at (-2.2,-2.199) {$+$};
  \node[anchor=base, inner sep=0pt] at (-2.2,-2.706) {$v_{gs}$};
  \node[anchor=base, inner sep=0pt] at (-2.2,-3.251) {$-$};
  \node[anchor=base, inner sep=0pt] at (-2.2,-3.949) {$S$};
  \node[anchor=base, inner sep=0pt, text=anot] at (1.0,-4.60) {$v_{gs} = -v_{id}/2$};

  % ---------------- legendas das partes ----------------
  % mesma linha de base para (a) e (b), longe o bastante da anotação para não formar legenda de duas linhas
  \node[anchor=base, inner sep=0pt] at (-7.6,-5.45) {(a) Meio-circuito diferencial};
  \node[anchor=base, inner sep=0pt] at (1.0,-5.45) {(b) Modelo $\pi$-híbrido};
\end{circuitikz}
\end{document}
